% TOP_PROBLEM: {"id": 13, "title": "ABC Conjecture", "category_id": 1, "category": {"id": 1, "name": "number_theory", "display_name": "Number Theory", "description": "Properties of integers, prime numbers, Diophantine equations.", "slug": "number-theory", "order_index": 1, "created_at": "2026-07-31T15:26:25.670Z"}, "status": "open"}
% FIRST_POSTED: null
% !TeX program = lualatex
% Compile twice: lualatex open_problems_500.tex
% All references and prose are embedded; no BibTeX database or external inputs.
\documentclass[11pt,a4paper]{article}
\usepackage[margin=24mm,headheight=14pt]{geometry}
\usepackage{fontspec}
\setmainfont{Latin Modern Roman}
\setsansfont{Latin Modern Sans}
\setmonofont{Latin Modern Mono}
\usepackage{amsmath,amssymb,mathtools}
\usepackage{unicode-math}
\setmathfont{Latin Modern Math}
\usepackage{microtype}
\usepackage{enumitem,xurl,needspace}
\usepackage[dvipsnames]{xcolor}
\usepackage{hyperref}
\hypersetup{unicode=true,colorlinks=true,linkcolor=MidnightBlue,urlcolor=MidnightBlue,
 pdftitle={Mathematical Research Notebook},pdfauthor={Source-annotated compilation},bookmarksnumbered=true}
\usepackage{bookmark}
\usepackage{tocloft}
\setlength{\cftsecnumwidth}{3em}
\cftsetpnumwidth{2.5em}
\cftsetrmarg{4em}
\usepackage{titlesec}
\titleformat{\section}{\Large\bfseries\raggedright}{\thesection}{0.7em}{}
\usepackage{fancyhdr}
\pagestyle{fancy}\fancyhf{}
\fancyhead[L]{\small\sffamily Open Problems Math | Research notebook}
\fancyhead[R]{\small 27 September 2026}
\fancyfoot[C]{\thepage}
\setlength{\parindent}{0pt}
\setlength{\parskip}{4pt plus 1pt}
\setlength{\emergencystretch}{3em}
\setcounter{tocdepth}{1}
\setcounter{secnumdepth}{1}
\setlist[itemize]{leftmargin=3em,itemsep=3pt,topsep=4pt,parsep=0pt}
\allowdisplaybreaks
\newcommand{\N}{\mathbb{N}}
\newcommand{\Z}{\mathbb{Z}}
\newcommand{\Q}{\mathbb{Q}}
\newcommand{\R}{\mathbb{R}}
\newcommand{\C}{\mathbb{C}}
\newcommand{\F}{\mathbb{F}}
\newcommand{\A}{\mathbb{A}}
\newcommand{\PP}{\mathbb P}
\newcommand{\E}{\mathbb{E}}
\newcommand{\eps}{\varepsilon}
\newcommand{\dd}{\,\mathrm d}
\newcommand{\sref}[2]{\hyperlink{source:#1:#2}{[S#2]}}
\newcommand{\eref}[2]{\hyperlink{extra:#1:#2}{[E#2]}}
% Turn the next switch off for a shorter reading copy. The quotations preserve
% every exactTarget field, including qualifications omitted from a short summary.
\newif\ifshowcatalogue
\showcataloguetrue
\newcommand{\cataloguescope}[1]{\ifshowcatalogue
 \paragraph{Exact scope in the supplied catalogue.}
 \begin{quote}\small #1\end{quote}\fi}
\usepackage{etoolbox}
\AtBeginEnvironment{itemize}{\small}
% Standalone entry; original section content is delimited for verification.
\begin{document}
\setcounter{section}{0}
%% BEGIN ORIGINAL SECTION CONTAINER
% ================================================================
% problemId: 13
% Canonical problem ID: 13
% exactTarget SHA-256: d3586c02f7e75ac8c61fca67676d74fe20e6e677853bd2f6e5fa2c94930ef5f8
\clearpage
\section*{\texorpdfstring{ABC Conjecture}{ABC Conjecture}}\label{13}
\subsection{Definitions and mathematical statement}
For a positive integer $n$, its radical is $\operatorname{rad}(n)=\prod_{p\mid n}p$, with $\operatorname{rad}(1)=1$. A triple $a,b,c>0$ with $a+b=c$ and $\gcd(a,b)=1$ is automatically pairwise coprime. The conjecture is
\[
 \forall\varepsilon>0\ \exists C_\varepsilon>0\ \forall a,b,c\in\Z_{>0},\qquad
 \begin{cases}a+b=c,\\\gcd(a,b)=1\end{cases}
 \Longrightarrow c<C_\varepsilon\operatorname{rad}(abc)^{1+\varepsilon}.
\]
The constant depends on $\varepsilon$ only. A counterexample must defeat every constant for some fixed positive $\varepsilon$; one unusually large triple does not suffice.
\par\smallskip\textit{Formulation and concept references:} \sref{13}{1}.
\cataloguescope{Prove that for every epsilon > 0, there exists a constant C(epsilon) such that for all coprime positive integers a + b = c, c < C(epsilon) rad(abc)\textasciicircum{}\{1+epsilon\}.}
\subsection{Short English statement}
Can the size of a coprime additive triple always be controlled, up to any small exponent loss, by its distinct prime factors?
%% BEGIN RESEARCH INSERTION
\subsection{Research attempt}
\paragraph{Current frontier.}
\textit{Research check: 27 September 2026.}
The selected assertion has exponent $1+\varepsilon$ for every
$\varepsilon>0$, not one fixed large exponent. Pasten's arithmetic-derivative
framework relates carefully quantified small-derivative statements to
$abc$-type bounds; it motivates the lattice attack below.
\eref{13}{1}, Sections 2--4. His September 2026 manuscript reports
improved discriminant--conductor bounds toward Szpiro, not the target
power bound here. \eref{13}{2}. Letendre's July 2026 paper proposes
another conjecture and conditional applications. \sref{13}{4}.
The inherited proof claims are not adopted as established results:
Scholze--Stix explain a specific objection to the IUT argument and record
that its proponents disagree. \sref{13}{3}. This notebook does not purport
to audit the complete IUT literature or Joshi's later claim.

\paragraph{Attack route.}
The polynomial $abc$ analogy suggests replacing differentiation by an
integer-valued map satisfying Leibniz and the single additive relation
$a+b=c$. The missing ingredient is then a small \emph{nondegenerate}
integer solution of a linear constraint. Rather than assume that any
short lattice vector works, we compute exactly the image of the
Wronskian functional on that lattice and test its transverse direction.
This is an investigation within an established framework, not a claim
to have invented arithmetic derivatives.

\paragraph{Attempt: the derivative lattice and its arithmetic image.}
Fix coprime positive $a+b=c$ with $(a,b)\ne(1,1)$, let
$P=\{p:p\mid abc\}$ and $R=\prod_{p\in P}p$. For integer weights
$w=(w_p)_{p\in P}$, extended by zero outside $P$, define
\[
 D_w(n)=n\sum_{p\mid n}\frac{v_p(n)}p w_p\quad(n\geq1).
\]
This is an integer and obeys $D_w(uv)=uD_w(v)+vD_w(u)$.
Define two integer row vectors indexed by $P$:
\[
 A_p=\begin{cases}a v_p(a)/p,&p\mid a,\\
 b v_p(b)/p,&p\mid b,\\-c v_p(c)/p,&p\mid c,
 \end{cases}\qquad
 B_p=\begin{cases}-ab v_p(a)/p,&p\mid a,\\
 ab v_p(b)/p,&p\mid b,\\0,&p\mid c.
 \end{cases} \tag{11.1}
\]
Pairwise coprimality makes these cases disjoint. The allowed weights
form $K=\ker(A:\Z^P\to\Z)$; for them
$D_w(a)+D_w(b)=D_w(c)$ and the Wronskian is
$W(w)=aD_w(b)-bD_w(a)=Bw$. The rows $A,B$ are independent over $\Q$:
$A$ has a nonzero coordinate at a prime dividing $c$, where $B$ vanishes,
and $B$ is nonzero at a prime dividing $ab$. Thus $B(K)\ne\{0\}$.

Here is an exact lattice lemma, proved rather than inferred from a norm
heuristic. Put $g=\gcd_p|A_p|$ and
\[
 d=\frac{\gcd_{p<q}|A_pB_q-A_qB_p|}{g}>0.
 \qquad B(K)=d\Z. \tag{11.2}
\]
To prove this, the primitive row $A/g$ can be transformed by a unimodular
integer column matrix $U$ to $(1,0,\ldots,0)$, using Bezout's identity
and successive Euclidean column operations. Write
$BU=(b_1,\ldots,b_s)$. In these coordinates $K$ is generated by the
last $s-1$ standard basis vectors, so its image is
$\gcd(b_2,\ldots,b_s)\Z$. The ideal generated by all two-by-two minors
is unchanged by unimodular column operations: each new minor is an
integer combination of old ones, and applying $U^{-1}$ gives the
reverse inclusion. For the transformed rows $A/g,B$, that ideal is
exactly $(b_2,\ldots,b_s)$. Multiplying the first row by $g$ proves
(11.2). In particular, the smallest nonzero \emph{Wronskian magnitude}
is explicitly computable from the valuations; this says nothing yet
about the height of weights achieving it.

The nonzero minors in (11.2) occur only between different prime blocks.
For example, if $p\mid a$ and $q\mid b$, the minor equals
$abc\,v_p(a)v_q(b)/(pq)$. Between the $a,c$ or $b,c$ blocks the same
formula holds up to sign with the corresponding valuations. This gives
a direct finite computation of $d$, without enumerating all $w$.

\paragraph{Attempt: divisibility versus height.}
For every $w\in K$, one has
\[
 \frac{abc}{R}\mid W(w). \tag{11.3}
\]
Indeed, if $p^e\Vert n$, the defining formula for $D_w(n)$ is divisible
by $p^{e-1}$: the $p$-term has this factor and every other term has
$p^e$. For $p\mid a$ or $p\mid b$, apply this to
$W=aD_w(b)-bD_w(a)$. For $p\mid c$, use the equal expression
$W=aD_w(c)-cD_w(a)$. Multiplying the resulting prime-power divisibilities
proves (11.3). This recovers the core divisibility in Pasten's
Theorem 3.3. \eref{13}{1}.

If $W(w)\ne0$ and $H=\|w\|_\infty$, then
\[
 \frac{abc}{R}\leq |W(w)|
 \leq abH\sum_{p\mid ab}\frac{v_p(ab)}p
 \leq \frac{abH\log(ab)}{2\log2}
 \leq \frac{abH\log c}{\log2}. \tag{11.4}
\]
The middle inequality follows termwise from $p\log p\geq2\log2$.
Thus $c\leq RH\log c/\log2$. The exact obstruction is now a transverse
minimum
\[
 \kappa(a,b)=\min\{\|w\|_\infty:w\in K,\ Bw\ne0\}, \tag{11.5}
\]
not the length of the shortest nonzero vector of $K$.
For example, $a=1,b=8,c=9$ gives $A=(12,-6)$, $B=(12,0)$,
$g=6$, $d=12$, and $w=(1,2)$ has $W=12=abc/R$.

A sufficient, deliberately stronger hypothesis would be: for every
$\eta>0$ there is $C_\eta$ such that
$\kappa(a,b)\leq C_\eta c^\eta$ for the nonexceptional triples described
below. Given a target $\varepsilon>0$, choose
$\eta>0$ and $\delta>0$ with $\eta+\delta<\varepsilon/(1+\varepsilon)$.
Using $\log c\leq C_\delta c^\delta$ in (11.4) yields
$c^{1-\eta-\delta}\ll R$, hence
$c\ll R^{1/(1-\eta-\delta)}\leq R^{1+\varepsilon}$.
Increasing the constant makes the final inequality strict. This supplies
the exact exponent bookkeeping, but the hypothesized transverse bound
is not proved or asserted plausible in this full strength.

\paragraph{Stress test: the short direction can be useless.}
For the abstract integer rows $A=(H,1,0)$ and $B=(1,0,0)$,
$K$ has the length-one vector $(0,0,1)$, yet every $w\in K$ with
$Bw\ne0$ has $\|w\|_\infty\geq H$. Moreover (11.2) gives $d=1$.
So even the conjunction of a very short kernel vector and a small
Wronskian image generator does not yield a short \emph{transverse}
vector. This is a countermodel to the lattice inference, not an
arithmetic counterexample to $abc$.

The arithmetic example $1+108=109$ gives
$A=(108,108,-1)$ and $B=(108,108,0)$ on primes $2,3,109$.
The vector $(1,-1,0)$ is short but has $W=0$; any nonzero $W$ forces
$|w_{109}|\geq108$. This illustrates the exceptional prime configurations
also discussed in \eref{13}{1}, Example 3.8. They should not be used to
claim an unconditional lower bound along an infinite sequence that has
not been constructed. Such triples are easy for $abc$ itself: if one
entry is $1$ and another is prime $q$, then $c\leq q+1\leq2R$.
They can therefore be excluded from the sufficient hypothesis after
handling them separately. The case $(1,1,2)$ is also immediate.

The remaining gap is not solved by computing a Smith form: that finds
integer generators and the exact image, but may produce very large
weights. Nor does a bound $\kappa\ll c^\eta$ for one fixed
$0<\eta<1$ alone give every exponent $1+\varepsilon$ through (11.4).
It gives only a fixed larger power after absorbing the logarithm. No
stronger implication from the literature is silently substituted for
this calculation.

\paragraph{Outcome.}
\textbf{Outcome: Exploratory approach with unresolved gap.}
The investigation gives an exact minors formula for all achievable
Wronskians, separates image size from the transverse minimum, and
exhibits both abstract and arithmetic failures of a shortest-vector
shortcut. The $abc$ exponent would follow from the stated much stronger
transverse estimate, but that estimate remains unproved. These are
not a proof of $abc$, a counterexample to it, or a claim of historical
novelty for the underlying derivative framework or lattice algebra.
%% END RESEARCH INSERTION
\subsection{Sources}
\begin{itemize}\raggedright
\item[\textnormal{[S1]}]\hypertarget{source:13:1}{} Distance Between Cubics and Rationals, Results in Mathematics (2026), statement of the abc conjecture.
\url{https://link.springer.com/article/10.1007/s00025-026-02616-5}.
{\small\textit{Catalogue formulation source.}}
\item[\textnormal{[S2]}]\hypertarget{source:13:2}{} Construction of Arithmetic Teichmuller Spaces IV: Proof of the abc-conjecture — Kirti Joshi.
\url{https://arxiv.org/abs/2403.10430}.
{\small\textit{Catalogue research/status reference; not a proof endorsement.}}
\item[\textnormal{[S3]}]\hypertarget{source:13:3}{} Why abc is still a conjecture — Peter Scholze and Jakob Stix, July 16, 2018.
\url{https://www.math.uni-bonn.de/people/scholze/WhyABCisStillaConjecture.pdf}.
{\small\textit{Catalogue research/status reference; not a proof endorsement.}}
\item[\textnormal{[S4]}]\hypertarget{source:13:4}{} The abc Conjecture Revisited — Patrick Letendre, July 2026.
\url{https://arxiv.org/abs/2607.07641}.
{\small\textit{Catalogue research/status reference; not a proof endorsement.}}
%% BEGIN RESEARCH REFERENCES
\item[\textnormal{[E1]}]\hypertarget{extra:13:1}{} Hector Pasten, \emph{Arithmetic derivatives through geometry of numbers}, arXiv:2106.16165v2 (20 July 2021), Sections 2--4, Theorem 3.3 and Example 3.8.
\url{https://arxiv.org/html/2106.16165v2}.
{\small\textit{Added research source. Arithmetic derivatives, Wronskians and the obstruction from small but degenerate derivatives. Consulted 27 September 2026.}}
\item[\textnormal{[E2]}]\hypertarget{extra:13:2}{} Hector Pasten, \emph{Improved bounds for Szpiro's conjecture}, arXiv:2609.17390v1 (15 September 2026), abstract and main bounds.
\url{https://arxiv.org/abs/2609.17390}.
{\small\textit{Added research source. Recent manuscript reporting discriminant--conductor estimates; not a proof audit or an abc resolution. Consulted 27 September 2026.}}
%% END RESEARCH REFERENCES
\end{itemize}

%% END ORIGINAL SECTION CONTAINER
\end{document}
